\documentclass{article}
\usepackage{graphicx}
\usepackage{xcolor}
\usepackage{soul}
\usepackage[utf8]{inputenc}
\usepackage{newunicodechar}
\usepackage{tikz}
\usetikzlibrary{positioning}
\usepackage[T1]{fontenc}
\usepackage{fontspec}
\usepackage{polyglossia}
\setmainlanguage{english}
\setotherlanguage{greek}
\usepackage{geometry}
\geometry{margin=1in}
\usepackage{csquotes}
\usepackage{enumitem}
\usetikzlibrary{arrows.meta, positioning, shapes.geometric}
\usepackage{fancyvrb}
\usepackage{setspace}
\newfontfamily\greekfont{Gentium Plus}
\newcommand{\gk}[1]{{\greekfont #1}}
\newunicodechar{₁}{$_1$}
\usepackage[backend=biber,style=authoryear]{biblatex}
\addbibresource{references.bib}
\usepackage[colorlinks=true, linkcolor=blue, urlcolor=blue, citecolor=blue]{hyperref}
\usepackage{changepage}
\usepackage{comment}
\newcommand{\ra}{\ensuremath{\rightarrow}\ }
\newcommand{\inlinenote}[1]{\par\noindent
  \fcolorbox{orange}{orange!20}{\parbox{0.97\linewidth}{#1}}\par\noindent}

% ============================================================================
% DRAFT-v1 — Part II Methodology section. Justification layer (M.0–M.4).
% Style target = Part I (Lanham fingerprint: 85% Germanic, voice 0.39, opacity
%   0.74, ~31w sentences, dynamicRange 0.90, heavy chiasmus/antithesis/polyptoton).
% Bib keys (calleja2011, burke1945/burke1950, rickert2013) need reconciling with
%   references.bib in the citation pass (Stage D). Secondary lit cited inline
%   (Author Page) / (\textit{ShortTitle} Page) per Part I convention.
% STILL TO COME: M.5 instruments+dwelling · M.6 VR↔screen differential ·
%   M.7 protocol · Appendix (full brawl walkthrough, v8-updated).
% ============================================================================

\begin{document}

\section{Rhetorical-Ontological Diachronic Analysis (RODA)}

The foregoing account discharged its proper burden and no more. It established that virtual and phantasmic things are real, that they are distinct in kind both from the non-virtual and from the grounds on which they depend, and that they share a single mode of being---a potency rendered into act; and, having fixed what such things \textit{are}, it deferred the question of how they come to work upon us, and even to reshape us, to the analysis that follows. That analysis is the method this section sets out. I call it Rhetorical-Ontological Diachronic Analysis (RODA): a method that reads a stretch of play as the \textit{actualization of desire}---the movement of a disposed, desiring being from a perceptual encounter, through a charged image and a committed judgment, to a completed action, and on into the lasting disposition (\textit{hexis}) that the action leaves behind.\footnote{Throughout, Burke's and Calleja's own claims are reported as theirs (their texts speak for themselves); every binding \textit{between} their frameworks, and every application to the behavior of a player, is the dissertation's own construction---an anticipatory reading neither author makes. The distinction matters for what follows: the method claims construction, never exegesis.} Where the metaphysics carried the reality-burden, the method carries the mechanism. It traces how a designed world drives that actualization toward a substantial union---what I will call \textit{rhetorical incorporation}---and how that union, once won, sediments into the player's dispositional ground and can, at the limit, reshape how one dwells in the world beyond the game.

The method is not, then, a frame borrowed from game studies or from rhetoric and afterward fitted to Aristotle. It is built from the actualization of desire---the ontology Part I constructs---and it takes up Calleja's \textit{incorporation} and Burke's \textit{consubstantiation} as matter for that form. What follows lays out, first, the ontological frame the method presupposes; then the generative chain and the union it drives toward; and then, at length, the two accounts the method re-grounds and fuses---for it is there, in the re-grounding, that the originality lies.

\subsection{The ontological frame, by reference}

The chain does not float. Every stretch of play the method reads is, in Aristotle's own terms, a motion of a substance, explained by the four causes; and the actualization-of-desire chain is the fine structure of one of those causes---the efficient-final causation of a desiring agent, the cell of the four-cause grid at which a rational animal's own motion is started and directed.\footnote{The frame is established in Part I (\S\S3--4) and is not re-derived here: the four causes (\textit{Met.} I.3, 983a24--32; \textit{Phys.} II.3, 194b24--195a3), the three-factor kinetic schema (\textit{DA} III.10, 433b12--26), and the seven causes of action (\textit{Rhet.} I.10, 1368b33--1369a6) that individuate the efficient-final complex for human action. I rehearse only what the method invokes downstream.} I will not re-lay that ground. But two features of it do specific work in what follows, and so must be set out here.

The first concerns how the four causes apply to a thing made to be used. Because the four causes belong to a single change, they cannot be parcelled out between two agents: designing and playing are two different actions, and each is owed its own complete four.\footnote{The corrective matters because the tempting error is to assign the player the efficient and final cause and the designer the material and formal of one and the same act. There is no one act there to be so divided. There are two acts---a making and a playing---joined at a hinge.} They are joined at Aristotle's relation of the producing-art to the using-art (\textit{Phys.} II.2, 194a36--b8): the art that produces ``knows the matter,'' the art that uses ``knows the form'' and is ``directive''; and because the form of a game is its being-played, the player's playing is the very for-the-sake-of-which the designer's making serves. One artifact, two angles---and the crossing this licenses is the spine of every later diagnosis: the designer's end-product is the player's given means. What the world \textit{supplies} (the material and formal cause of the playing) is furnished by the designer's making; what the player \textit{does} (the efficient and final cause) remains their own. The designer cannot be the efficient cause of the player's act---one agent cannot supply another's desire---but only alters their material world to solicit a response. ``Why now'' is the designer's doing; ``which way'' is the player's.

The second feature is the keystone the union will inherit: Aristotle's account, in \textit{De Anima} III.7, of a single actuality that admits two accounts---``one and the same actualization, though their being is not the same'' (431a12--13). It is on this structure, carried over from perception to the union of agent and world, that the central event of the method turns.

\subsection{The generative spine and the union it drives toward}

The chain organizes Aristotle's analyses of perception and action into five nodes joined by four motions. At $A_0$ the world stands open as a horizon---a sensible object and an organ of sense, given together within motion and time, the designed \textit{Umwelt} into which the player is already thrown. The perceptual motion completes at $A_1$, where a perception is reached and a dual residue laid down---an \textit{eidetic} trace and a hedonic-\textit{orectic} one. The residue settles, at $A_2$, into the \textit{phantasma} proper, the charged image on which every later operation works. At $A_3$ the image is met by the soul's cognitive powers and ratified by \textit{doxa} as how things stand, and the determinate form desire takes---the emotion---concretizes. And at $A_4$ the action completes, converging the prior causes into bodily movement and sedimenting, as it goes, into \textit{hexis}. The chain is read forward, motion by motion; but, because each completed actuality is what makes the preceding motion intelligible, it is understood backward, from the object of desire (\textit{orekton}) that draws the whole.

At the transition from $A_3$ to $A_4$ the chain reaches a substantial union, and the union has two faces held as one event. On the world-face the player becomes one with the world they act \textit{in}; on the other-face they become one with the avatar they act \textit{as} and the others the player acts \textit{with}---substantially one, yet a distinct locus, both joined and separate. This is the method's original core, and the claim is precise: neither account the method inherits has \textit{both} faces, and neither grounds the union where the method grounds it---in the actualization of desire. I name the fused event \textit{rhetorical incorporation}, and the name drives the work the fusion attempts---\textit{incorporation} (the world-face) made \textit{consubstantial} (the other-face) and rendered ontological by the chain that generates them. The union is an \textit{experiential} convergence; the player becomes one \textit{with} the world and the others. This does not a collapse the difference in kind I established in my previous analysis of the metaphysics of virtual and phantasmic things. A player incorporated to the hilt has not thereby made the virtual non-virtual; the union joins the player to the world they act in, it does not dissolve the difference in kind between them.\footnote{This guard is the same two-in-being clause the keystone already carries (\textit{DA} III.7, 431a12--13): the union was always ``one substrate, \textit{two in being}.'' Its development as a disruption-axis instrument---the foreclosed rupture, the \textit{Gestell} limit---and the media differential that turns on it are taken up below (\S\S M.5--M.6) and demonstrated on their proper beats in the analysis that follows.}

\subsection{The world-face: incorporation, re-grounded (Calleja)}

The chain tells us that a designed world presses on the efficient-final cause of the player's act. It does not, on its own, tell us \textit{through what} the world presses, nor how one might know---from what a player does, and says she undergoes---that the union of player and world was reached at all. For these the dissertation turns to Gordon Calleja, and it turns to him because no account built closer to the philosophy of mind could supply what his, built from games, does.

Calleja's \textit{In-Game: From Immersion to Incorporation} sets aside ``immersion'' and ``presence''---the metaphors of transport, of a subject plunged into a container-world---and offers in their place \textit{incorporation}: a union that runs in both directions at once, in which ``the environment is assimilated'' into the player's awareness \textit{and} the player is, by way of the avatar, embodied within that environment, the two held together rather than taken in turn (\textit{In-Game} 169). The model Calleja develops he calls the Player Involvement Model (PIM): six dimensions of involvement---kinesthetic, spatial, shared, narrative, affective, ludic---that Calleja draws inductively from years of play and from interviews conducted in-game, as clusters of emphasis rather than posited kinds (\textit{In-Game} 36, 43--44). What the chain lacks from within itself---a validated, fine-grained, games-built vocabulary for where a designed world bears on the soul's motion, and a test for when the world-face union obtains---Calleja already has.

The method I put forward here preserves this, transforms it, and lifts it into the fusion. It \textit{preserves} incorporation as the world-face, and the six dimensions as the designed channels through which a world applies its pressure to the chain---named, in an analysis, only where the world is in fact pressing, never run as an inventory to fill.\footnote{The recasting of the six dimensions as ``channels'' through which the designed efficient cause presses on the chain is the dissertation's construction, not Calleja's; he offers them as dimensions of involvement, and the binding to the chain's motions is mine.} It \textit{transforms} the concept's ground. For Calleja incorporation is a psychological intensification of involvement---``the absorption of a virtual environment into consciousness''---a formula that runs, finally, from world to mind; on the reading I develop, it is neither an absorption of the world by a mind nor a mind's projection into the world, but the player-substance and the world-substance becoming one substrate, two in being (\textit{DA} III.7): the player \textit{dwelling} in the designed \textit{Umwelt}, gathered by it as much as gathering it. And it \textit{lifts} the world-face into union with the other-face, which Calleja's model, world-facing throughout, never reaches.

That the re-grounding is warranted, rather than imposed, Calleja's own text invites. He grants that incorporation is ``ultimately a metaphor'' and his model ``a first considered step'' (\textit{In-Game} 5); and his official ground---an experientialist, embodied-cognition naturalism---sits uneasily beside the phenomenological vocabulary he in fact reaches for, the ``sense of habitation,'' the ``inhabiting'' of a place, the continuity of the played world with the lived one. The vocabulary outruns the naturalism that is meant to underwrite it. What the method offers is not a correction of Calleja against his will but the ground his own commitments leave unfilled: the two halves interlock---Calleja gives the method its games-built flesh, and the actualization of desire gives that flesh its ontological skeleton. Neither, alone, does the work.

This is also where Calleja earns his place as an instrument and not an ornament. His own writing supplies a test for incorporation as against mere involvement---a small set of conditions on which, fail any one, the state falls back to involvement---and the method takes up that test, numbers and extends it as a differentiator (R1--R7), and makes it the gauge of the union's world-face: assimilation and embodiment, held \textit{at once} (the strictest gatekeeper), on a spatial-kinesthetic cornerstone, several dimensions blended, deepening with duration, and porous to the real world rather than sealed against it.\footnote{R3--R7 are Calleja's own five-component criterion (\textit{In-Game} 169--173, 178); the numbering and the two added conditions (R1 assimilation, R2 embodiment, drawn from his double-axis) are the method's operationalization. The full differentiator and its provenance are given in Appendix B.} The union-claim is thereby kept from assertion: it is read off conditions, and it returns a verdict---achieved, partial, merely involved, or ruptured.

It is worth saying why Calleja, and not another. The rival accounts of game experience either keep the very dualism the method means to dissolve or lack the grain to locate the world's pressure. Murray's transport---immersion, agency, transformation---is the unidirectional ``plunge'' Calleja himself absorbs and corrects, a subject lowered into a world, with no dimensional decomposition and no in-game evidence beneath it; Ryan's immersion stays narratological and textual, rich on the spatial axis Calleja keeps as his cornerstone but silent on the kinesthetic, the shared, the ludic; the formalist tradition carries the magic-circle and the ``immersive fallacy'' that Calleja's continuity-thesis is built to excise. The decisive thing is the non-dualism. Where the rivals keep a subject set over against an object-world, Calleja's double-axis already crosses that divide---the world taken into awareness as the player is given over into the world---and it is just this crossing that maps onto the keystone's overarching one in substance and different in being formulation, and lets incorporation become the world-face of a single event rather than one more name for a subject's absorption.

\subsection{The other-face and the rhetoric of motive (Burke)}

Incorporation joins the player to the world they act in. It does not join them to the gang they act with; and the chain, which tells us \textit{what} moved the player, does not tell us how that motion is afterward rendered, nor how its true source comes to be concealed. Two residues are left, then---a union of agents that can itself \textit{source} an act, and a grammar of the motive an act is given---and for both the dissertation turns to Kenneth Burke.

Take the first. Aristotle's \textit{praxis} is owned by a single soul, and \textit{thymos}, the spirited desire, ``needs a toward-whom'' that Aristotle never raises to a union of agents; Calleja's incorporation is player and world, and social only as co-presence, not as a substance that motivates. Burke alone supplies the union from which a collective desire issues. To act together, on his account, is to be made consubstantial: two agents become ``substantially one,'' yet each ``an individual locus of motives,'' both joined and separate (\textit{Rhetoric} 21); and, the union is compensatory to a prior division (\textit{Rhetoric} 22)---a division the player arrives in more literally than any reader, her own body set over against the avatar they plays \textit{as}. Because consubstantiation is built of acting-together, it can source the next act. In the Valentine saloon the point is plain: Arthur does not wade into the brawl from a private grievance but because he is one with the gang---the shove that lands on Bill lands, consubstantially, on him. The keystone is thereby shown to be generative, not decorative: the union does not merely crown the act, it sources it.

Take the second. The pentad---act, scene, agent, agency, purpose, and the later sixth term, attitude---is not, in the method's hands, a grid of boxes to fill. It is read as Burke read it, as a grammar of \textit{motive-attribution}: the terms and their ratios name what a finished act gets ascribed to, and the ratios let one and the same act be relocated from motive to motive---onto the scene (``the situation made me''), onto the tools, onto the others---so that the term most often left unspoken is the one that did the work, the player's own sedimented \textit{hexis}.\footnote{The pentad is used only as Burke uses it, diagnostically and not as a classification to be exhausted; the rule of use is that a term is named where it illuminates a stage of the chain or the union, never run as a five- or six-box pass. Attitude is Burke's own sixth term (his 1962 addendum), and the method reads it as a \textit{hexis}---the disposition-with-which an act is done.} Where the chain gives what moved the agent, Burke gives how the agent, and the analyst, accounts for the moving---and how the account conceals its own ground. RODA \textit{preserves} the term consubstantiation; it \textit{transforms} it from a bond of symbolic persuasion into the desire-sourced ``WE'' that \textit{doxa} ratifies at $A_3$ and the joint act consummates at $A_4$; and it \textit{lifts} it into the other-face of rhetorical incorporation. 

Why Burke, and not the nearer alternatives? Aristotle's own \textit{Rhetorica} details the emotions and the means of persuasion, but it has no consubstantial union of agents and no pentadic grammar of motive; to press it into this service would be to make it reinvent precisely what Burke already provides. Perelman's adherence of minds and Bitzer's rhetorical situation remain on the side of an audience persuaded---propositional, scenic---and offer no substance made one, no act that the union sources. What sets Burke apart is that his consubstantiality is drawn from his own paradox of \textit{substance}, and so docks onto Aristotle's ontology of act and potency without importing an alien epistemology: Burke himself reads Aristotle on the act and on \textit{entelechy}---the priority in form by which ``man is `prior' to boy'' (\textit{Grammar} 261)---a reading the dissertation has already leaned upon. Among the theorists of motive he is the one whose framework is itself an account of action, and the one whose apparatus, set beside Calleja's, independently converges with it.

That convergence is the last and strongest reason the two belong together. When the pentad and the PIM are set side by side and matched, term against dimension, they meet---blind, having been built for wholly different ends---on a small shared core, and, at the center of it, on the single binding of consubstantiation to incorporation.\footnote{The bidirectional crosswalk and its four-binding core (scene--spatial, agent--shared, agency--kinesthetic, attitude--affective) are retained as Appendix A, where they function now as evidence rather than as method: two pipelines that converge blind are tracking one structure.} That two such different instruments should come to rest at the same place is evidence that there is one structure beneath the two descriptions; and it is that one structure the method names rhetorical incorporation. The world-face and the other-face are not two lenses laid over the same scene. They are the two faces the convergence brings to light.

% ============================================================================
% TO COME (next drafting pass):
%   M.5 — the method's emergent instruments + dwelling (compressed; Group-A
%         instruments DERIVED + forward-referenced per decision D6: firewall,
%         foreclosed-rupture/Gestell, existence-fragility).
%   M.6 — the VR↔screen differential module (concluding module; route-not-
%         reachability; data-discrepancy axis as the metaphysics→method seam;
%         phantasia-load the linchpin; G&G illustrative gesture; Sonny posed
%         as OPEN question).
%   M.7 — the protocol in brief + scale variants (worksheet → appendix pointer).
%   Appendix A — the demoted Burke↔Calleja crosswalk (evidence, not method).
%   Appendix B — the R1–R7 incorporation differentiator (provenance).
%   Appendix C — the full RODA walkthrough of the saloon brawl (the advisor
%         demonstration, updated to v8: prohairesis = bouletic's completion,
%         NOT a coordinate Type 4; "phantasmic" not "phantasmatic"; no bold
%         run-in heads; restored Step 3 = the Burke motive-rhetoric step).
% ============================================================================

\end{document}
